
Mākslīgie neironu tīkli ir kļuvuši par nozīmīgu tehnoloģiju datorredzes un autonomo sistēmu jomā, dažos gadījumos sasniedzot
vai pat pārsniedzot cilvēka smadzeņu veiktspēju. Tomēr to darbības pamatā ir matricu reizināšanas skaitliski aprēķini, un tie patērē
ievērojami vairāk enerģijas nekā bioloģiskās neironu sistēmas. Tas apgrūtina to izmantošanu iegultajās skaitļošanas
platformās, īpaši reāllaika apstrādes uzdevumos, kuros tiek izmantoti reti notikumu sensoru dati. Neiromorfā
skaitļošana, kuras pamatā ir pulsējoši neironu tīkli (SNN), ir alternatīva pieeja, kas cenšas
atdarināt bioloģisko neironu uz notikumiem balstīto komunikāciju un aprēķinus.

Lai gan neiromorfie algoritmi un notikumu sensori, piemēram, dinamiskās redzes un kohleārie sensori, ir ievērojami
attīstījušies, aparatūras paātrināta skaitļošana joprojām ir aktuāls izaicinājums. Lielākā daļa dziļo SNN modeļu joprojām
tiek novērtēti, izmantojot programmatūras simulācijas ietvarus uz tradicionālās CPU un GPU aparatūras. Tādēļ SNN potenciālās
energoefektivitātes priekšrocības bieži vien ir ierobežotas.

Specializēti aparatūras paātrinātāji, piemēram, ASIC un FPGA, ir spējīgi novērst šos ierobežojumus, izpildot retos
aprēķinus specializētā aparatūrā. Īpaši FPGA nodrošina līdzsvaru starp elastību un energoefektivitāti, ļaujot izstrādāt un
pētīt specializētus paātrinātājus, izvairoties no sarežģītām un dārgām pielāgotu ASIC risinājumu implementācijām.

Šajā darbā tiek pētīta dziļo pulsējošu neironu tīklu uz FPGA balstīta paātrināšana attēlu atpazīšanas uzdevumiem. Darbs
pievēršas uz pārkonfigurējamas aparatūras arhitektūras izstrādi, pulsējošu neironu modeļu implementāciju un pulsējošu
konvolucionālo neironu tīklu praktisku pielietojumu 2D un 3D uztveres uzdevumos.

Izstrādātais paātrinātājs implementē notikumu konvolūciju, konfigurējamus noplūstošus integrē-un-pulsē (LIF) un Ižikeviča
neironu modeļus, kā arī adrešu-notikumu reprezentācijas (AER) komunikāciju pārkonfigurējamā, uz PCIe balstītā FPGA ietvarā.
% Lai demonstrētu arhitektūras pielietojamību reālu problēmu risināšanā, tā tiek novērtēta attēlu klasifikācijas un 3D objektu
% detektēšanas uzdevumos, izmantojot notikumu attēlu, LiDAR un lidojuma laika punktu mākoņu datu kopas.
FPGA realizācija ir eksperimentāli novērtēta, izmantojot attēlu klasifikācijas uzdevumus. Atsevišķi ir izstrādāta
un novērtēta 3D-SCNN modeļa arhitektūra objektu noteikšanai, izmantojot LiDAR un lidojuma laika (ToF) sensoru punktu mākoņu datu
kopas, demonstrējot impulsu neironu tīklu praktisku pielietojamību telpiskās vides uztveres uzdevumos.

Eksperimentālie rezultāti parāda, ka FPGA implementācija nodrošina reāllaika inferenci ar ievērojami mazāku enerģijas
patēriņu salīdzinājumā ar GPU platformām, vienlaikus saglabājot zemu latentumu un augstu caurlaidību. Turklāt izstrādātie
pulsējošie konvolucionālie neironu tīkli ir pielāgojami arī punktu mākoņu sensoru datu modalitātei, uzsverot notikumu
uztveres pielāgojamību automatizētām iegultajām sistēmām, kas paļaujas uz punktu mākoņiem balstītu uztveri.
